\documentclass[aspectratio=169]{beamer}
\usetheme{Madrid}
\usecolortheme{default}
\usepackage{booktabs}
\usepackage{graphicx}
\usepackage{amsmath}

% Clean, professional look
\setbeamertemplate{navigation symbols}{}
\setbeamertemplate{footline}[frame number]

\title{How Accurate Are Learned World Models?}
\subtitle{Midterm Progress Report --- DreamerV3 Imagination Fidelity on Snake}
\author{Josh Manchester}
\institute{CS 5080 --- Reinforcement Learning \\ University of Colorado Colorado Springs}
\date{March 31, 2026}

\begin{document}

\begin{frame}
\titlepage
\end{frame}

%------------------------------------------------------------
\begin{frame}{Recap: Research Question}

\textbf{Central Question:} Does DreamerV3 need accurate imagination to learn effective policies, or can it succeed even when its world model is wrong?

\vspace{1em}
\textbf{Three Hypotheses:}
\begin{itemize}
    \item[\textbf{H1:}] Imagination accuracy improves as training progresses
    \item[\textbf{H2:}] Better imagination leads to better play
    \item[\textbf{H3:}] Semantic accuracy matters more than pixel accuracy
\end{itemize}

\vspace{1em}
\textit{Why it matters:} DreamerV3 trains its policy \textbf{entirely} inside imagined trajectories. If those trajectories are wrong, the agent is optimizing for scenarios that do not exist.
\end{frame}

%------------------------------------------------------------
\begin{frame}{Recap: Why Snake?}

\begin{columns}
\begin{column}{0.55\textwidth}
\textbf{Snake as testbed:}
\begin{itemize}
    \item $10 \times 10$ grid, $64 \times 64$ RGB observations
    \item 4 discrete actions, deterministic transitions
    \item Reward: $+1$ food, $-1$ collision, $-0.002$/step
    \item Distinct colors $\rightarrow$ automated state extraction
\end{itemize}

\vspace{0.5em}
\textbf{Why it works:}\\
When the model gets something wrong, we know \textit{exactly} what it got wrong. No visual clutter, no ambiguity.
\end{column}
\begin{column}{0.42\textwidth}
\textbf{Two levels of measurement:}\\[0.5em]
\textbf{Pixel:} MSE, SSIM\\
(``Do the images look alike?'')\\[0.5em]
\textbf{Semantic:} Head error, body accuracy, food correct\\
(``Does the model know the game state?'')
\end{column}
\end{columns}
\end{frame}

%------------------------------------------------------------
\begin{frame}{What Has Been Completed}

\textbf{Five training configurations --- all complete with full analysis:}

\begin{table}
\centering
\small
\begin{tabular}{llccl}
\toprule
\textbf{Config} & \textbf{Obs} & \textbf{Steps} & \textbf{LR} & \textbf{Key Change} \\
\midrule
Baseline    & 64$\times$64 & 1.1M & 1e-4 & --- \\
Extended    & 64$\times$64 & 2M   & 1e-4 & Fine-tuned, longer episodes \\
Slow LR     & 64$\times$64 & 2M   & 3e-5 & $\frac{1}{3}\times$ learning rate \\
Res-32      & 32$\times$32 & 530K & 1e-4 & Lower resolution \\
Res-16      & 16$\times$16 & 500K & 1e-4 & Lowest resolution \\
\bottomrule
\end{tabular}
\end{table}

\vspace{0.5em}
Also trained a $3\times$ LR model --- world model learned the game, but policy never stabilized.\\
Model 6 (doubled train ratio) currently at 801K / 2M steps.
\end{frame}

%------------------------------------------------------------
\begin{frame}{Game Performance Across Configurations}

\begin{table}
\centering
\begin{tabular}{lccc}
\toprule
\textbf{Config} & \textbf{Mean Return} & \textbf{Peak} & \textbf{Final 5 Avg} \\
\midrule
Extended    & 31.17 & 54.26 & 36.50 \\
Slow LR     & 16.57 & 24.60 & 19.16 \\
Baseline    &  9.82 & 25.71 & 22.28 \\
Res-32      &  1.06 & 13.60 &  9.83 \\
Res-16      &  0.78 &  9.13 &  7.44 \\
\bottomrule
\end{tabular}
\end{table}

\vspace{0.5em}
\begin{itemize}
    \item Resolution \textbf{dramatically} affects gameplay
    \item Fine-tuning + longer episodes gives the biggest boost
    \item Slow LR competitive but still improving at 2M steps
\end{itemize}
\end{frame}

%------------------------------------------------------------
\begin{frame}{Preliminary H1: Imagination Improves Over Training}

\begin{columns}
\begin{column}{0.45\textwidth}
\begin{table}
\centering
\small
\begin{tabular}{lcc}
\toprule
\textbf{Config} & \textbf{Impr.} & \textbf{Sig.} \\
\midrule
Res-16      & 5/5 & 5/5 \\
Res-32      & 5/5 & 5/5 \\
Baseline    & 5/5 & 4/5 \\
Slow LR     & 5/5 & 2/5 \\
Extended    & 3/5 & 0/5 \\
\bottomrule
\end{tabular}
\end{table}
\end{column}
\begin{column}{0.52\textwidth}
\textbf{Emerging pattern:}
\begin{itemize}
    \item H1 strongest in \textbf{actively-learning} models
    \item Extended model inherited a good world model $\rightarrow$ flat imagination trends
    \item World model learns visuals \textbf{quickly}, then plateaus
    \item Policy learning is the \textbf{slow} part
\end{itemize}
\end{column}
\end{columns}

\vspace{0.5em}
Baseline head error: 0.268 $\rightarrow$ 0.023 grid cells ($\rho = -0.85$, $p < 0.001$)\\[0.3em]
\textbf{Assessment:} H1 is well-supported. Expect this to hold in remaining runs.
\end{frame}

%------------------------------------------------------------
\begin{frame}{Preliminary H2: Imagination Correlates with Performance}

\begin{table}
\centering
\small
\begin{tabular}{lccccc}
\toprule
\textbf{Metric} & \textbf{Base} & \textbf{Ext} & \textbf{Slow} & \textbf{R32} & \textbf{R16} \\
\midrule
MSE              & .570$^{*}$  & .099 & .023 & .603$^{*}$  & .564$^{*}$  \\
SSIM             & .630$^{**}$ & .120 & .066 & .660$^{**}$ & .641$^{*}$  \\
Head err         & .843$^{***}$& .163 & .514$^{*}$  & .513$^{*}$  & .678$^{**}$ \\
Body acc         & .352        & .511$^{*}$  & .740$^{***}$& .292 & .532$^{*}$  \\
Food corr        & .597$^{**}$ & .090 & .310 & .654$^{**}$ & .636$^{**}$ \\
\bottomrule
\end{tabular}
\end{table}

\textbf{Look at the Slow LR column:} Pixel metrics $\approx$ 0. Semantic metrics strong.\\[0.3em]
\textbf{Assessment:} H2 is well-supported. Head error and body accuracy are reliable predictors. Extended model is weaker because the bottleneck shifted from world model to policy.
\end{frame}

%------------------------------------------------------------
\begin{frame}{Preliminary H3: Semantic vs. Pixel --- The Key Finding}

\begin{table}
\centering
\begin{tabular}{lcccc}
\toprule
\textbf{Config} & \textbf{Pixel} $\overline{|\rho|}$ & \textbf{Semantic} $\overline{|\rho|}$ & \textbf{Fisher $z$} & $p$ \\
\midrule
\textbf{Slow LR} & \textbf{0.045} & \textbf{0.522} & \textbf{2.65} & \textbf{0.008} \\
Extended    & 0.110 & 0.255 & ---  & n.s. \\
Baseline    & 0.600 & 0.597 & 1.42 & 0.154 \\
Res-32      & 0.631 & 0.486 & ---  & n.s. \\
Res-16      & 0.603 & 0.615 & ---  & n.s. \\
\bottomrule
\end{tabular}
\end{table}

\vspace{0.3em}
Slow LR: Semantic metrics predict performance \textbf{$>$10$\times$} better than pixel metrics.

\vspace{0.3em}
\textbf{Why only Slow LR?} The reduced learning rate \textbf{decouples} visual reconstruction from semantic learning. Under normal LR, both improve together, masking the distinction.

\vspace{0.3em}
\textbf{Assessment:} One significant result so far. Need remaining runs to determine if this generalizes or is specific to the slow LR regime.
\end{frame}

%------------------------------------------------------------
\begin{frame}{Resolution Experiment Insights}

\textbf{Same game, different visual fidelity:}

\vspace{0.5em}
\begin{columns}
\begin{column}{0.48\textwidth}
\textbf{Lower resolution produces:}
\begin{itemize}
    \item Worse gameplay (peak 9 vs 54)
    \item But cleaner statistical signals (5/5 significant H1 vs 0/5)
    \item World model achieves high accuracy quickly on simple visual task
\end{itemize}
\end{column}
\begin{column}{0.48\textwidth}
\textbf{Preliminary interpretation:}
\begin{itemize}
    \item Good dreams about a blurry world are not enough
    \item Agent needs spatial detail for \textit{planning}
    \item Confirms analysis pipeline detects real patterns
    \item Useful control for the final paper
\end{itemize}
\end{column}
\end{columns}
\end{frame}

%------------------------------------------------------------
\begin{frame}{Remaining Work and Timeline}

\textbf{In progress:}
\begin{itemize}
    \item Model 6 (doubled train ratio): 801K / 2M steps --- estimated complete April 3
    \item Will test whether more imagination per real step improves learning
\end{itemize}

\vspace{0.5em}
\textbf{Planned for final paper:}
\begin{itemize}
    \item Model 7: $20 \times 20$ grid, $128 \times 128$ obs --- does the world model scale?
    \item Collision prediction accuracy (new semantic metric)
    \item Longer imagination horizons --- how fast do predictions degrade?
    \item Cross-configuration synthesis of all H1/H2/H3 results
\end{itemize}

\vspace{0.5em}
\textbf{Key question for final paper:} Does the H3 result ($p = 0.008$) replicate under other training conditions, or is it specific to the slow LR regime?

\vspace{0.5em}
\textbf{Hardware:} NVIDIA RTX 5070 Ti (16 GB), PyTorch 2.7.0, CUDA 12.8
\end{frame}

%------------------------------------------------------------
\begin{frame}{Summary --- Where Things Stand}

\begin{enumerate}
    \item \textbf{H1:} Well-supported. Imagination improves over training, strongest in actively-learning models.
    \item \textbf{H2:} Well-supported. Head error and body accuracy are the most reliable predictors ($\rho$ up to $-0.84$).
    \item \textbf{H3:} Significant in one configuration ($p = 0.008$). Semantic $>$ pixel by $>$10$\times$ under slow LR. Remaining runs will test generalizability.
\end{enumerate}

\vspace{0.5em}
\textbf{Emerging big picture:} DreamerV3 does not need visually accurate dreams. It needs dreams that capture the right \textit{structure}. The world model's grasp of game semantics---not pixel quality---appears to drive policy learning.

\vspace{0.5em}
\textbf{On track for final paper} --- analysis pipeline operational, 5/7 planned runs complete.

\end{frame}

%------------------------------------------------------------
\begin{frame}{}
\centering
\Huge Questions?
\end{frame}

\end{document}
